\documentclass[10pt,twocolumn,a4paper]{article}
\usepackage[margin=1.7cm,columnsep=0.65cm]{geometry}
\usepackage{times}
\usepackage{graphicx}
\usepackage{booktabs}
\usepackage{multirow}
\usepackage{amsmath,amssymb}
\usepackage{siunitx}
\usepackage{microtype}
\usepackage{enumitem}
\usepackage{url}
\usepackage[hidelinks]{hyperref}
\usepackage{xcolor}
\usepackage{caption}
\usepackage{subcaption}
\usepackage{float}
\usepackage{array}
\usepackage{tabularx}
\usepackage{cite}

\definecolor{notegray}{RGB}{90,90,90}
\newcommand{\isian}[1]{\textcolor{notegray}{\textit{[Isi: #1]}}}
\newcommand{\projecttitle}{Judul Final Project Pengolahan Citra Digital dan Visi Komputer}

\title{\bfseries \projecttitle}
\author{
Nama Mahasiswa 1$^{1}$, Nama Mahasiswa 2$^{1}$, Nama Mahasiswa 3$^{1}$\\[-1mm]
\small $^{1}$Program Studi ..., Fakultas Ilmu Komputer, Universitas Brawijaya\\
\small Email: mahasiswa1@student.ub.ac.id, mahasiswa2@student.ub.ac.id, mahasiswa3@student.ub.ac.id
}
\date{}

\begin{document}
\maketitle

\begin{abstract}
\isian{150--250 kata. Jelaskan masalah nyata di lingkungan Universitas Brawijaya, data primer yang dikumpulkan, tahapan PCD utama, metode Machine Learning, metode Deep Learning/Computer Vision, skenario perbandingan, metrik evaluasi, hasil utama, dan kesimpulan. Jangan menulis sitasi di abstrak.}
\end{abstract}

\noindent\textbf{Keywords:} \isian{3--5 kata kunci, misalnya image processing, computer vision, machine learning, deep learning, primary dataset.}

\section{Pendahuluan}
\isian{Jelaskan latar belakang masalah dan manfaatnya. Masalah harus relevan dengan objek/kondisi yang dapat diamati di lingkungan UB. Jelaskan mengapa pendekatan pengolahan citra dan visi komputer diperlukan.}

\isian{Tuliskan ringkasan penelitian/metode terkait dari referensi ilmiah yang relevan. Tunjukkan celah atau alasan project ini tetap perlu dilakukan. Hindari sekadar daftar penelitian.}

\isian{Paragraf terakhir harus memuat tujuan dan kontribusi project. Contoh kontribusi: (1) dataset primer yang diambil sendiri di lingkungan UB; (2) pipeline PCD; (3) perbandingan PCD+ML, Raw+DL, dan PCD+DL; (4) evaluasi kuantitatif dan analisis error.}

\section{Data Primer dan Akuisisi Data}
\subsection{Objek dan Lokasi Akuisisi}
\isian{Jelaskan objek, kelas/target, lokasi pengambilan data di dalam lingkungan Universitas Brawijaya, waktu/kondisi pengambilan, dan alasan objek dipilih. Data wajib primer dan dikumpulkan sendiri. Hindari objek sensitif/identitas personal sebagai target utama.}

\subsection{Prosedur Akuisisi}
\isian{Jelaskan perangkat kamera, resolusi, jarak, sudut, pencahayaan, latar, jumlah data per kelas, serta variasi kondisi. Jelaskan bagaimana kelompok menjaga konsistensi sekaligus variasi dataset.}

\begin{table}[H]
\centering
\caption{Ringkasan dataset primer.}
\label{tab:dataset}
\begin{tabular}{lrrr}
\toprule
Kelas & Train & Validation & Test\\
\midrule
Kelas A & -- & -- & --\\
Kelas B & -- & -- & --\\
Kelas C & -- & -- & --\\
\bottomrule
\end{tabular}
\end{table}

\subsection{Pelabelan dan Pembagian Data}
\isian{Jelaskan proses labeling dan pembagian train/validation/test. Pastikan tidak terjadi data leakage. Jika data berupa video, frame yang sangat berdekatan tidak boleh tersebar secara sembarang ke train dan test.}

\section{Metodologi}
\subsection{Gambaran Umum Sistem}
\isian{Tampilkan diagram pipeline lengkap: Data Primer UB $\rightarrow$ PCD $\rightarrow$ fitur/ROI $\rightarrow$ ML/DL $\rightarrow$ evaluasi $\rightarrow$ demo.}

\begin{figure}[H]
\centering
\fbox{\parbox[c][3.0cm][c]{0.92\columnwidth}{\centering Tempatkan diagram pipeline project di sini}}
\caption{Pipeline keseluruhan sistem.}
\label{fig:pipeline}
\end{figure}

\subsection{Pengolahan Citra Digital}
\isian{Jelaskan tahapan PCD yang benar-benar digunakan dan alasan teknisnya, misalnya resize, RGB/HSV/LAB, grayscale, gamma, CLAHE, mean/Gaussian/median/bilateral, sharpening, DFT/FFT, restorasi, Canny, Otsu/adaptive threshold, morfologi, atau CCL. Jangan memasukkan metode hanya untuk memperbanyak tahapan.}

Contoh formulasi yang dapat digunakan bila relevan:
\begin{equation}
 s = 255\left(\frac{r}{255}\right)^{\gamma},
 \label{eq:gamma}
\end{equation}
\isian{Definisikan seluruh simbol pada persamaan dan jelaskan pengaruh parameter $\gamma$. Ganti/tambah persamaan sesuai metode yang digunakan.}

\subsection{Ekstraksi Fitur dan Machine Learning}
\isian{Jelaskan fitur handcrafted yang digunakan, misalnya histogram warna, GLCM, LBP, HOG, Hu Moments, contour/shape, SIFT/ORB, dan lain-lain. Jelaskan minimal dua metode ML klasik yang dibandingkan, misalnya SVM, Random Forest, atau KNN. Sertakan parameter penting.}

\subsection{Deep Learning / Computer Vision}
\isian{Jelaskan model DL/CV yang digunakan, misalnya CNN, MobileNet, YOLO, U-Net, Mask R-CNN, atau model lain yang sesuai. Jelaskan input, training/transfer learning, augmentasi, hyperparameter utama, dan output model.}

\subsection{Skenario Eksperimen}
Tiga eksperimen minimum yang \textbf{wajib} dibandingkan adalah:
\begin{enumerate}[leftmargin=*]
    \item \textbf{PCD + Handcrafted Features + Machine Learning};
    \item \textbf{Raw Image + Deep Learning};
    \item \textbf{PCD + Deep Learning}.
\end{enumerate}
\isian{Tambahkan ablation study bila memungkinkan, misalnya Raw $\rightarrow$ +CLAHE $\rightarrow$ +Filtering $\rightarrow$ +Segmentation, untuk menunjukkan dampak setiap tahapan PCD.}

\section{Hasil dan Pembahasan}
\subsection{Hasil Pengolahan Citra}
\isian{Tampilkan citra sebelum dan sesudah PCD. Jangan hanya menyatakan ``lebih baik''; jelaskan perubahan yang terjadi dan hubungkan dengan kebutuhan model.}

\begin{figure}[H]
\centering
\fbox{\parbox[c][3.0cm][c]{0.92\columnwidth}{\centering Contoh: Raw | CLAHE | Filtering | Segmentasi}}
\caption{Contoh hasil tahapan pengolahan citra.}
\label{fig:pcdresult}
\end{figure}

\subsection{Hasil Machine Learning dan Deep Learning}
\begin{table}[H]
\centering
\caption{Perbandingan hasil eksperimen utama.}
\label{tab:result}
\resizebox{\columnwidth}{!}{%
\begin{tabular}{lcccc}
\toprule
Metode & Accuracy & Precision & Recall & F1\\
\midrule
PCD + ML-1 & -- & -- & -- & --\\
PCD + ML-2 & -- & -- & -- & --\\
Raw + DL & -- & -- & -- & --\\
PCD + DL & -- & -- & -- & --\\
\bottomrule
\end{tabular}}
\end{table}
\isian{Sesuaikan metrik dengan tugas. Object detection: Precision, Recall, mAP@0.5, mAP@0.5:0.95. Segmentation: IoU dan Dice. Counting: MAE/RMSE. Tracking: metrik tracking yang relevan.}

\subsection{Analisis Pengaruh PCD}
Salah satu bentuk analisis sederhana adalah:
\begin{equation}
\Delta F1 = F1_{\mathrm{PCD+Model}} - F1_{\mathrm{Raw+Model}}.
\end{equation}
\isian{Analisis apakah preprocessing meningkatkan atau justru menurunkan performa. Jelaskan penyebabnya berdasarkan karakteristik data, bukan hanya angka.}

\subsection{Analisis Keberhasilan dan Kegagalan}
\isian{Tampilkan beberapa contoh prediksi benar dan salah. Analisis penyebab error: pencahayaan, blur, occlusion, background, objek kecil, kemiripan antarkelas, segmentasi gagal, dan sebagainya.}

\subsection{Efisiensi dan Kelayakan Implementasi}
\isian{Jika relevan, laporkan waktu inferensi, ukuran model, kebutuhan komputasi, atau FPS. Bahas apakah sistem realistis untuk digunakan pada skenario yang diusulkan.}

\section{Demo / Prototipe Sistem}
\isian{Jelaskan bentuk demo final: input citra/video, proses utama, output prediksi/deteksi/segmentasi/tracking, serta cara pengguna menjalankan sistem. Sertakan screenshot antarmuka bila ada.}

\section{Pembagian Kontribusi Anggota}
\begin{table}[H]
\centering
\caption{Kontribusi anggota kelompok.}
\label{tab:contribution}
\resizebox{\columnwidth}{!}{%
\begin{tabular}{p{2.4cm}p{5.2cm}}
\toprule
Anggota & Kontribusi Utama\\
\midrule
Mahasiswa 1 & Data primer, EDA, dan Pengolahan Citra Digital.\\
Mahasiswa 2 & Ekstraksi fitur, Machine Learning, dan evaluasi ML.\\
Mahasiswa 3 & Deep Learning/Computer Vision, integrasi, dan demo.\\
\bottomrule
\end{tabular}}
\end{table}
\isian{Tuliskan kontribusi aktual dan spesifik. Akuisisi data tetap harus melibatkan seluruh anggota. Semua anggota wajib memahami keseluruhan pipeline.}

\section{Kesimpulan}
\isian{Jawab tujuan project secara ringkas berdasarkan hasil. Nyatakan metode yang diuji, dampak PCD, perbandingan ML dan DL, keterbatasan, serta satu arah pengembangan yang realistis. Jangan memasukkan hasil baru.}

\section*{Pernyataan Dataset Primer}
Kami menyatakan bahwa dataset utama pada project ini merupakan \textbf{data primer yang dikumpulkan sendiri oleh kelompok di dalam lingkungan Universitas Brawijaya}. \isian{Tambahkan ringkasan lokasi dan periode akuisisi.}

\section*{Ketersediaan Kode dan Data}
\isian{Tuliskan tautan repository/folder pengumpulan yang ditentukan dosen. Jangan membuka data yang mengandung informasi sensitif.}

\bibliographystyle{IEEEtran}
\bibliography{references}

\end{document}
